\section{Revisión de literatura}
La productividad agregada puede mejorar porque las empresas producen mejor, porque las más productivas ganan peso o porque cambia el conjunto de productores. La revisión compara qué parte de esa historia identifica cada método y qué encuentran sus aplicaciones. Reúne veinte referencias formales; no todas aplican Melitz--Polanec ni todas utilizan la misma medida de productividad.

La descomposición distingue cuatro fuentes: el cambio dentro de las empresas que permanecen, la reasignación entre ellas y las diferencias de productividad asociadas a entradas y salidas. La literatura ayuda a entender qué significa cada aporte y por qué sus resultados varían entre países.

\subsection{Qué mide cada componente}
\textbf{Cambio dentro de las empresas.} Compara la productividad media de las empresas presentes al inicio y al final, dando el mismo peso a cada una. Un aporte positivo indica que esa media aumentó; uno negativo, que disminuyó.

\textbf{Reasignación entre empresas.} Indica si las empresas más productivas aumentaron su peso relativo frente a las menos productivas. Un aporte positivo señala que la asociación entre productividad y participación se fortaleció; uno negativo, que se debilitó. Olley y Pakes (1996) establecieron esta separación entre productividad media y composición; Melitz y Polanec (2015) la extendieron para incorporar entradas y salidas.

\textbf{Entrada y salida.} La entrada compara a las empresas nuevas en la muestra con las que permanecen al final. La salida compara a las que permanecen con las que dejan de observarse, usando su productividad al inicio. Así, una entrada negativa indica que las entrantes tenían menor productividad que las sobrevivientes; una salida positiva, que las salientes tenían menor productividad que quienes permanecieron. Griliches y Regev (1995), Foster, Haltiwanger y Krizan (2001) y Melitz y Polanec (2015) muestran que la magnitud de estos aportes depende del método y del horizonte.

\subsection{La evidencia cambia entre países y periodos}
En Colombia, Eslava et al. (2004, 2006) encuentran que la reasignación y la salida de plantas menos productivas contribuyeron a aumentos de productividad tras las reformas de los años noventa. Para un periodo más reciente, el Banco Mundial (2024) atribuye un papel importante al cambio dentro de las empresas manufactureras; el FMI (2024), usando PTF y Orbis, encuentra que la reasignación amortigua la debilidad interna en manufactura y agricultura.

En Chile, el diagnóstico de IFC (2022) encuentra para 2010--2015 una media negativa, reasignación positiva, entrada negativa y salida positiva. Los mismos signos aparecen en nuestra EAM para ese periodo. El parecido sugiere una pregunta común: cómo puede aumentar el agregado cuando la productividad media de las empresas existentes cae. En ambos casos la composición contribuye a explicarlo, aunque Chile mide PTF y pondera por ventas, mientras nuestro ejercicio mide productividad laboral y pondera por empleo.

México ofrece un contraste. Iacovone et al. (2021) encuentran que los cambios entre empresas sobrevivientes dominan y que la reasignación reduce la productividad agregada durante buena parte de 1993--2018. En servicios de Portugal, el FMI (2022) documenta una entrada negativa asociada a diferencias de tamaño: las nuevas microempresas pueden elevar el promedio dentro de su grupo y, aun así, estar por debajo de las pymes sobrevivientes. Estos casos ayudan a leer los signos de Colombia sin asumir que sean universales.

Otros estudios refuerzan la importancia del contexto. En Israel, Griliches y Regev (1995) encuentran que predomina el cambio dentro de las empresas. En el Reino Unido, Harris y Moffat (2019) muestran que la reasignación compensó una media negativa en la manufactura. Ben Hassine (2019) encuentra que la contribución atribuida a la entrada varía según se use MP, GR o FHK. Song (2021) documenta una contribución positiva de la salida de plantas menos productivas en Corea.

\subsection{Qué aporta esta evidencia al caso colombiano}
La literatura permite interpretar los resultados del boletín en tres pasos. Primero, el resultado industrial —aumento agregado con cambio interno negativo— también aparece en Chile y el Reino Unido. Segundo, la evidencia de México e Israel recuerda que la reasignación puede ser negativa o pequeña; no es una ley general que las empresas productivas ganen peso. Tercero, el contraste entre Banco Mundial (2024) y nuestro resultado industrial exige revisar la unidad de observación, el periodo y el procesamiento antes de proponer una explicación económica.

Las comparaciones numéricas requieren igualar la medida de productividad, los pesos, la unidad observada, el periodo y la cobertura. Los resultados de estudios con PTF no son directamente equivalentes a los de productividad laboral. Además, la contribución contable describe el cambio observado; por sí sola no identifica sus causas.

\textbf{Fuentes y alcance.} Esta revisión reúne veinte referencias formales. Algunas se consultaron mediante sus textos completos y otras mediante resúmenes o secciones específicas; las fichas de Literatura detallan ese alcance. Brown et al. (2018) se usa solo como contexto regional por acceso parcial. La revisión es dirigida y puede ampliarse en las siguientes etapas.
\clearpage
